\section{MODEL\_PREDICTIVE/ems\_offline\_mpc\_v0.py}

\subsection*{Scopo}
Runner offline MPC: prepara dati reali e forecast, esegue controllo MPC step-by-step e salva risultati.

\subsection*{CLI}
\begin{itemize}[leftmargin=1.2cm]
\item \texttt{--config}: YAML MPC.
\item \texttt{--MPC\_inference\_dataset\_path}: dataset usato dal simulatore (opzionale se presente in config).
\item \texttt{--MPC\_forecast\_dataset\_path}: dataset usato dal controller MPC (opzionale, fallback a config e poi a inference dataset).
\item \texttt{--start-step}, \texttt{--end-step}: finestra simulata.
\end{itemize}

\subsection*{Flusso}
\begin{itemize}[leftmargin=1.2cm]
\item Carica configurazione e seleziona la sorgente dati:
\texttt{scenario.dataset\_mode=pymgrid\_bundle} usa i tre file scenario (\texttt{load/pv/grid}),
altrimenti usa il CSV \texttt{inference} tradizionale.
\item In modalita \texttt{pymgrid\_bundle} i prezzi buy/sell sono letti dal dataset \texttt{GridModule};
in modalita CSV senza colonne prezzo, i prezzi sono derivati da \texttt{price\_bands}.
\item Risolve i path dataset con priorita:
CLI \textrightarrow{} \texttt{ems.*}/\texttt{scenario.*} nel config \textrightarrow{} fallback forecast=inference.
\item Costruisce microgrid offline.
\item Istanzia \texttt{MPCController} (da \texttt{mpc\_MILP.py}).
\item A ogni step invoca \texttt{get\_action()} e applica \texttt{microgrid.step()}.
\item Salva log in cartella run timestampata sotto \texttt{outputs/MPC}.
\item Il \texttt{final\_report.txt} include il campo \texttt{Degradazione SoH : ON/OFF} ricavato dalla configurazione batteria.
\item Estrae cronologia transizioni batteria solo se disponibile (eventuale manager/transition model); con batteria base \texttt{pymgrid} la cronologia puo essere assente.
\item Il conteggio overshoot e letto in modo opzionale (fallback a \texttt{0} se l'attributo non e presente).
\item Mostra tabelle finali con \texttt{pandasgui} usando dataframe normalizzati (colonne flat); il \texttt{forecast\_log} viene mostrato solo se contiene dati.
\end{itemize}

\subsection*{Note manutenzione}
Qualsiasi cambio a nomi argomenti CLI va propagato in script test/sweep e README.
